\documentclass[aspectratio=169,11pt]{beamer}
% Shared CMPSC 480 Beamer preamble.
\usepackage[T1]{fontenc}
\usepackage[utf8]{inputenc}
\usepackage{lmodern}
\usepackage{amsmath,amssymb,mathtools}
\usepackage{booktabs,array,tabularx}
\usepackage{xcolor}
\usepackage{listings}
\usepackage{tikz}
\usetikzlibrary{arrows.meta,positioning,shapes.geometric}

% Ignore only negligible TeX box diagnostics that are below visual tolerance.
% Larger layout defects still appear in the build log and in rendered QA.
\vfuzz=3pt
\hbadness=2000

\definecolor{courseink}{HTML}{111111}
\definecolor{coursegray}{HTML}{EDEDED}
\definecolor{courserule}{HTML}{B8BCC4}
\definecolor{courseblue}{HTML}{3D8DFF}
\definecolor{courselightblue}{HTML}{D0EDFA}
\definecolor{coursemuted}{HTML}{5C6470}
\definecolor{coursegreen}{HTML}{0B7A53}
\definecolor{coursered}{HTML}{B3261E}

\usetheme{default}
\usefonttheme{professionalfonts}
\setbeamertemplate{navigation symbols}{}
\setbeamersize{text margin left=0.65cm,text margin right=0.65cm}

\setbeamercolor{normal text}{fg=courseink,bg=white}
\setbeamercolor{frametitle}{fg=courseink,bg=white}
\setbeamercolor{title}{fg=courseink,bg=white}
\setbeamercolor{subtitle}{fg=coursemuted,bg=white}
\setbeamercolor{structure}{fg=courseblue}
\setbeamercolor{alerted text}{fg=coursered}
\setbeamercolor{block title}{fg=courseink,bg=courselightblue}
\setbeamercolor{block body}{fg=courseink,bg=coursegray}
\setbeamercolor{example text}{fg=coursegreen}

\setbeamerfont{title}{size=\fontsize{25}{29}\selectfont,series=\bfseries}
\setbeamerfont{subtitle}{size=\fontsize{13}{16}\selectfont}
\setbeamerfont{frametitle}{size=\fontsize{19}{22}\selectfont,series=\bfseries}
\setbeamerfont{normal text}{size=\fontsize{11}{14}\selectfont}
\setbeamerfont{footline}{size=\fontsize{7.5}{9}\selectfont}

\setbeamertemplate{frametitle}{%
  \nointerlineskip
  % Use content-driven height so a long assessment title can wrap without
  % being clipped by a fixed-height title bar.
  \begin{beamercolorbox}[wd=\paperwidth,sep=0.17cm,leftskip=0.48cm,rightskip=0.48cm]{frametitle}
    \insertframetitle
  \end{beamercolorbox}
  \vspace{-0.08cm}
  {\color{courserule}\rule{\textwidth}{0.35pt}}
}

\setbeamertemplate{footline}{%
  \leavevmode
  \begin{beamercolorbox}[wd=\paperwidth,ht=2.6ex,dp=1.1ex,leftskip=.65cm,rightskip=.65cm]{author in head/foot}
    \color{coursemuted}\insertshorttitle\hfill\insertframenumber/\inserttotalframenumber
  \end{beamercolorbox}
}

\setbeamertemplate{itemize item}{\color{courseblue}\small$\blacktriangleright$}
\setbeamertemplate{itemize subitem}{\color{courseblue}\scriptsize$\bullet$}
\setbeamertemplate{enumerate item}{\color{courseblue}\insertenumlabel.}

\lstdefinestyle{coursepython}{
  language=Python,
  basicstyle=\ttfamily\fontsize{8.5}{10}\selectfont,
  keywordstyle=\color{courseblue}\bfseries,
  commentstyle=\color{coursemuted}\itshape,
  stringstyle=\color{coursegreen},
  showstringspaces=false,
  columns=fullflexible,
  keepspaces=true,
  frame=single,
  rulecolor=\color{courserule},
  backgroundcolor=\color{coursegray},
  breaklines=true,
  tabsize=4,
  xleftmargin=0.15cm,
  xrightmargin=0.15cm,
  framexleftmargin=0.15cm,
  framexrightmargin=0.15cm
}
\lstset{style=coursepython}

\newcommand{\points}[1]{\hfill{\color{courseblue}\textbf{[#1 points]}}}
\newcommand{\deliverable}[1]{\textbf{\color{courseblue}#1}}
\newcommand{\code}[1]{\texttt{#1}}
\newcommand{\keyidea}[1]{%
  \begin{beamercolorbox}[rounded=false,sep=0.55em,wd=\linewidth]{block body}
    \textbf{Key idea.} #1
  \end{beamercolorbox}
}
\newcommand{\answerlabel}{\textbf{\color{coursegreen}Answer.}\ }
\newcommand{\gradinglabel}{\textbf{\color{courseblue}Grading guidance.}\ }

\newenvironment{questionblock}[1]
  {\begin{block}{#1}}
  {\end{block}}

\newenvironment{solutionblock}[1]
  {\begin{block}{#1}}
  {\end{block}}

\AtBeginSection[]{
  \begin{frame}
    \vfill
    {\usebeamerfont{title}\color{courseink}\insertsectionhead\par}
    \vspace{0.4cm}
    {\color{courseblue}\rule{0.32\paperwidth}{3pt}}
    \vfill
  \end{frame}
}


% CMPSC480_TOTAL_POINTS: 150
% CMPSC480_ITEM_IDS: 1,2,3,4,5,6,7,8,9,10
\title[Final Examination Questions]{CMPSC 480 Comprehensive Final Examination}
\subtitle{Weeks 1-15: methods, evidence, integration, and responsibility}
\author{CMPSC 480 - Student Question Deck}
\date{}

\begin{document}


\begin{frame}
  \titlepage
\end{frame}

\begin{frame}{Assessment overview}
\begin{columns}[T,onlytextwidth]
\begin{column}{0.62\textwidth}
\Large Solve applied, mathematical, algorithmic, and engineering problems under explicit assumptions.

\vspace{0.35cm}
\normalsize
Coverage: All modules with emphasis on machine learning, evaluation, fairness, deployment, and cross-module integration
\end{column}
\begin{column}{0.33\textwidth}
\begin{block}{Points}150 total\end{block}
\begin{block}{Expected time}180 minutes\end{block}
\begin{block}{Submission}Individual examination PDF\end{block}
\end{column}
\end{columns}
\end{frame}

\begin{frame}{Learning objectives}
\begin{itemize}
  \item Compare and apply classical and adversarial search.
  \item Reason about stochastic planning and reinforcement learning.
  \item Perform Bayesian exact-inference operations.
  \item Derive regression, classification, and backpropagation updates.
  \item Design valid evaluation and hyperparameter protocols.
  \item Analyze fairness and responsible deployment decisions.
  \item Integrate methods through testable engineering contracts.
\end{itemize}
\end{frame}

\begin{frame}{Requirements checklist}
\small
\begin{itemize}
  \item Work independently and show intermediate algebra, frontier state, or factor state.
  \item One double-sided handwritten letter-size reference sheet is permitted.
  \item A nonprogrammable calculator is permitted; networked devices and communication are prohibited.
  \item State every tie-breaking, terminal, reward, and zero-division convention you use.
  \item Answers receive credit for justified reasoning; an unsupported final number may receive little credit.
  \item Academic integrity rules of the course and institution apply throughout the examination.
\end{itemize}
\end{frame}

\begin{frame}{Task 1 - Search and heuristic guarantees (15 points)}
\small
A* runs on a positive-cost graph with a heuristic that is admissible but not consistent.

\vspace{0.2cm}
\begin{enumerate}
\item[\textbf{a.}] State admissibility and consistency formally. \hfill{\color{courseblue}\textbf{[5 points]}}
\item[\textbf{b.}] Can graph-search A* remain optimal with this heuristic? State the needed implementation behavior. \hfill{\color{courseblue}\textbf{[5 points]}}
\item[\textbf{c.}] Design a test that exposes missing reopening. \hfill{\color{courseblue}\textbf{[5 points]}}
\end{enumerate}
\end{frame}

\begin{frame}{Task 2 - Adversarial reasoning (15 points)}
\small
A stochastic game node chooses action A with outcomes 8 and 0 at probabilities 0.25 and 0.75, or action B with guaranteed value 3.

\vspace{0.2cm}
\begin{enumerate}
\item[\textbf{a.}] Compute expectimax values and choose an action. \hfill{\color{courseblue}\textbf{[5 points]}}
\item[\textbf{b.}] What would a worst-case MIN interpretation choose for A, and why is that a different model? \hfill{\color{courseblue}\textbf{[5 points]}}
\item[\textbf{c.}] Give two tests for a combined minimax-expectimax implementation. \hfill{\color{courseblue}\textbf{[5 points]}}
\end{enumerate}
\end{frame}

\begin{frame}{Task 3 - MDPs and reinforcement learning (15 points)}
\small
For Q(s,a)=2, reward=-1, gamma=0.9, max-next-Q=5, and alpha=0.1, compare one Q-learning update with a Bellman optimality backup.

\vspace{0.2cm}
\begin{enumerate}
\item[\textbf{a.}] Compute the Q-learning target, TD error, and updated Q. \hfill{\color{courseblue}\textbf{[5 points]}}
\item[\textbf{b.}] How would the terminal case differ? \hfill{\color{courseblue}\textbf{[5 points]}}
\item[\textbf{c.}] Compare information required by value iteration and Q-learning. \hfill{\color{courseblue}\textbf{[5 points]}}
\end{enumerate}
\end{frame}

\begin{frame}{Task 4 - Bayesian inference (15 points)}
\small
A disease prior is 0.02, test sensitivity is 0.9, and false-positive rate is 0.1.

\vspace{0.2cm}
\begin{enumerate}
\item[\textbf{a.}] Compute P(disease|positive). \hfill{\color{courseblue}\textbf{[5 points]}}
\item[\textbf{b.}] Explain how variable elimination avoids repeated enumeration. \hfill{\color{courseblue}\textbf{[5 points]}}
\item[\textbf{c.}] Compare rejection sampling with likelihood weighting for rare evidence. \hfill{\color{courseblue}\textbf{[5 points]}}
\end{enumerate}
\end{frame}

\begin{frame}{Task 5 - Regression and gradient descent (15 points)}
\small
X=[[1,0],[1,1],[1,2]], y=[1,3,5], and w=[0,0] for one-half mean squared error.

\vspace{0.2cm}
\begin{enumerate}
\item[\textbf{a.}] Compute the initial residual and gradient. \hfill{\color{courseblue}\textbf{[5 points]}}
\item[\textbf{b.}] With alpha=0.1, compute the new weights. \hfill{\color{courseblue}\textbf{[5 points]}}
\item[\textbf{c.}] Explain how feature scaling and collinearity affect optimization or solving. \hfill{\color{courseblue}\textbf{[5 points]}}
\end{enumerate}
\end{frame}

\begin{frame}{Task 6 - Neural networks and backpropagation (15 points)}
\small
A layer has A\_prev shape (4,20), W shape (3,4), b shape (3,1), and delta shape (3,20).

\vspace{0.2cm}
\begin{enumerate}
\item[\textbf{a.}] Give Z, A, dW, and db shapes. \hfill{\color{courseblue}\textbf{[5 points]}}
\item[\textbf{b.}] Write the previous-layer delta formula and shape. \hfill{\color{courseblue}\textbf{[5 points]}}
\item[\textbf{c.}] Describe a two-sided gradient check and one limitation. \hfill{\color{courseblue}\textbf{[5 points]}}
\end{enumerate}
\end{frame}

\begin{frame}{Task 7 - Classification and regularization (15 points)}
\small
A binary classifier has logit z=log(4), label y=1, w=[3,4], lambda=0.5, and m=5.

\vspace{0.2cm}
\begin{enumerate}
\item[\textbf{a.}] Compute sigmoid(z) and cross-entropy loss. \hfill{\color{courseblue}\textbf{[5 points]}}
\item[\textbf{b.}] Compute the L2 penalty and gradient addition. \hfill{\color{courseblue}\textbf{[5 points]}}
\item[\textbf{c.}] Explain why threshold 0.5 is not universally optimal. \hfill{\color{courseblue}\textbf{[5 points]}}
\end{enumerate}
\end{frame}

\begin{frame}{Task 8 - Evaluation and tuning (15 points)}
\small
A rare-event classifier yields TP=12, FP=8, FN=3, and TN=977.

\vspace{0.2cm}
\begin{enumerate}
\item[\textbf{a.}] Compute accuracy, precision, recall, and F1. \hfill{\color{courseblue}\textbf{[5 points]}}
\item[\textbf{b.}] Training and validation error are both high and close. Diagnose and name one experiment. \hfill{\color{courseblue}\textbf{[5 points]}}
\item[\textbf{c.}] Give a valid hyperparameter workflow with k-fold CV. \hfill{\color{courseblue}\textbf{[5 points]}}
\end{enumerate}
\end{frame}

\begin{frame}{Task 9 - Fairness and responsible deployment (15 points)}
\small
Group A0 has TPR=0.85 and FPR=0.10; group A1 has TPR=0.65 and FPR=0.25. A production feature distribution then shifts.

\vspace{0.2cm}
\begin{enumerate}
\item[\textbf{a.}] Compute the equalized-odds gaps and interpret them. \hfill{\color{courseblue}\textbf{[5 points]}}
\item[\textbf{b.}] Explain why removing the protected attribute alone does not establish fairness. \hfill{\color{courseblue}\textbf{[5 points]}}
\item[\textbf{c.}] Specify a complete shift monitor and response. \hfill{\color{courseblue}\textbf{[5 points]}}
\end{enumerate}
\end{frame}

\begin{frame}{Task 10 - Cross-module integration (15 points)}
\small
Design a system that combines a learned heuristic, Bayesian belief state, and a search or control policy for navigation under uncertain hazards.

\vspace{0.2cm}
\begin{enumerate}
\item[\textbf{a.}] Define the component handoffs. \hfill{\color{courseblue}\textbf{[5 points]}}
\item[\textbf{b.}] State three component or end-to-end tests. \hfill{\color{courseblue}\textbf{[5 points]}}
\item[\textbf{c.}] Design a baseline, metrics, and limitation statement. \hfill{\color{courseblue}\textbf{[5 points]}}
\end{enumerate}
\end{frame}

\begin{frame}{Edge cases and defensive programming}
\small
\begin{itemize}
  \item A logistic probability is numerically zero.
  \item A subgroup metric denominator is zero.
  \item A ReLU check lands at zero.
  \item A learned heuristic overestimates.
  \item A production dependency fails during evaluation.
\end{itemize}
\end{frame}

\begin{frame}{Submission instructions}
\small
\begin{itemize}
  \item Write your name and identify every requested problem and subpart.
  \item Show enough intermediate work that another reader can audit the result.
  \item Submit the complete question-response document when time is called.
\end{itemize}
\vspace{0.2cm}
\alert{Do not submit credentials, generated caches, or unapproved libraries.}
\end{frame}

\begin{frame}{Grading rubric - part 1}
\scriptsize
\begin{tabularx}{\textwidth}{>{\bfseries}p{0.28\textwidth} p{0.08\textwidth} X}
\toprule
Category & Points & Required evidence \\
\midrule
Search and heuristic guarantees & 15 & Guarantee analysis, implementation rule, and counterexample reasoning. \\
Adversarial reasoning & 15 & Expectimax calculation, minimax contrast, and engineering test. \\
MDPs and reinforcement learning & 15 & Numeric update and planning-learning comparison. \\
Bayesian inference & 15 & Posterior, factor reasoning, and approximate-inference choice. \\
Regression and gradient descent & 15 & Gradient calculation, update, and conditioning analysis. \\
\bottomrule
\end{tabularx}
\end{frame}

\begin{frame}{Grading rubric - part 2}
\scriptsize
\begin{tabularx}{\textwidth}{>{\bfseries}p{0.28\textwidth} p{0.08\textwidth} X}
\toprule
Category & Points & Required evidence \\
\midrule
Neural networks and backpropagation & 15 & Shape derivation, gradient formulas, and checking. \\
Classification and regularization & 15 & Probability, loss, L2 penalty, and threshold interpretation. \\
Evaluation and tuning & 15 & Metrics, curve diagnosis, and selection protocol. \\
Fairness and responsible deployment & 15 & Disparity, criterion reasoning, and monitoring control. \\
Cross-module integration & 15 & Architecture, verification, and responsible evidence plan. \\
\bottomrule
\end{tabularx}
\end{frame}

\begin{frame}{Reflection prompt}
In one or two sentences:

Name one connection among modules that changed how you would design or evaluate an AI system.
\end{frame}

\end{document}
